% Fichier complété par tools/complete_missing_corrections.py.
% Source : COR/COR-03-PI/70_Hublex/70_Hublex.tex
% Statut : rédigé (9 question(s)).
\begin{corrigebox}[Corrigé rédigé]
\CorrigeQuestion{1}
Pour conserver l'unité et obtenir une comparaison directe courant
            consigne/courant mesuré, on choisit $K_{iu}=1$ dans les unités normalisées (ou
            l'inverse du gain du capteur si les unités électriques ne sont pas normalisées).
\CorrigeQuestion{2}
Pour un échelon unitaire, la valeur finale donne
            $\mu_s=1/[1+K_{iu}K_pK_m]$ pour une correction proportionnelle.
\CorrigeQuestion{3}
On compare cette expression aux tolérances de 1.7.1.1. Une erreur nulle ne
            peut être obtenue avec un gain fini sans intégrateur.
\CorrigeQuestion{4}
Le PI augmente fortement le gain basse fréquence et annule l'erreur
            statique, tandis que son zéro permet de préserver la phase autour de la
            coupure. Il accroît toutefois la sensibilité au bruit et peut réduire la marge.
\CorrigeQuestion{5}
Pour $C(p)=K_p+K_i/p=K_p(1+\omega_i/p)$,
            $\omega_i=K_i/K_p=100$ rad/s. Le gain haute fréquence vaut 20 dB et la phase
            passe de $-90^\circ$ à $0^\circ$ autour de $100$ rad/s.
\CorrigeQuestion{6}
On place d'abord le zéro environ une décade sous la pulsation de coupure
            visée, puis on translate le gain afin d'obtenir la marge imposée :
            $K_p=1/|L_0(j\omega_c)C_0(j\omega_c)|$.
\CorrigeQuestion{7}
Une fois $K_p$ fixé, $K_i=K_p\omega_i$.
\CorrigeQuestion{8}
Le réglage peut satisfaire la précision mais conduire à une commande trop
            forte ou à une bande passante incompatible avec l'actionneur et le bruit; il
            faut donc saturation et anti-windup.
\CorrigeQuestion{9}
Le constructeur ajoute typiquement une limitation de courant/rampe et un
            anti-windup de l'intégrateur, voire un préfiltre de consigne, afin de respecter
            les contraintes physiques et le dépassement.
\end{corrigebox}
